% ===== sec/10_discussion.tex =====
\section{Discussion}
\label{sec:discussion}

We discuss five matters: the meaning of a Type-0 abstraction on a
bounded laptop, the role of the start language, the producer--consumer
restatement and the subtlety of a Type-3 producer, threats to
validity, and the relationship between the trace-language theory
and the prior literature on agent architectures.

\subsection{Type-0 abstraction versus bounded execution}

The Goal-Driven producer is analysed through an idealised
recursively enumerable envelope.  In practice, the producer runs on
a single laptop with a finite tape and a finite execution budget,
supplemented by the High Performance Computing Center at Texas Tech
University for the heavier ablation runs.  The realised execution is
therefore bounded.  Lemma~\ref{lem:bounded-soundness} supports the
one-way comparison from bounded executions to the idealised envelope;
it does not prove that the physical artifact realizes a strict
Type-0 language.

The asymmetry matters.  A lower-bound claim on the producer's
expressive power would require a separate language-theoretic proof,
and a single observed trace cannot supply one.  We therefore use the
Type-0 abstraction as an envelope for recursive enumeration, not as a
strict empirical classification.  The Type-3 consumer brings the
narrow conformance-checking problem inside the decidable; it does
not decide membership in the full producer language.

In Hartmanis--Stearns'
language~\cite{hartmanis1965computational}, the realised machine
inhabits a complexity class strictly inside Type-1.  In
Turing's~\cite{turing1936computable} language, it is a finite
restriction of a universal machine.  Both views are correct;
neither is more correct than the other.

\subsection{Why the start language is necessary}

It would be tempting to dispense with the start language and ask
the Goal-Driven producer to invent its own operation alphabet.
We attempted this in pilot runs and found, predictably, that the
producer's enumeration drifts into pathological corners of
$\Sigma^*$ that the consumer was not designed to cover.  The
start language $L_0$ serves two formal roles:
\begin{enumerate}
\item It fixes the seed alphabet $\Sigma_0$ over which the
consumer's keyword set $K$ is defined.  Without $\Sigma_0$, the
Aho--Corasick DFA cannot be constructed.
\item It provides a finite syntactic skeleton from which the
Type-0 enumeration can extend in a structured way.  The producer
treats $L_0$ as a base case for its recursion, not as a fixed
ceiling.
\end{enumerate}
The start language is therefore the inductive hypothesis of the
recursive enumeration; without it, the recursion has no base.
The empirical CSV \texttt{trace\_language\_original.csv} is the
artefact form of the inductive hypothesis.  The realised CSV
\texttt{trace\_language\_kaggle.csv} is the conclusion of one
inductive step.  More inductive steps (one per task instance)
would produce more emergent columns; the architecture
generalises straightforwardly to the full \textsc{MLE-bench}
suite.

\subsection{Producer--consumer, and the Type-3 producer subtlety}
\label{sec:type3-producer}

The Koohestani--Li--Podkopaev--Izadi
question~\cite{koohestani2025automata} asks \emph{what class is
this agent?}\ and stops there.  The producer--consumer
question of Dijkstra~\cite{dijkstra1968cooperating} asks
\emph{how do classes \textbf{interact}?}\ --- which is the
question this paper answers.  In the canonical setting, the
producer runs the harder enumeration and the consumer runs the
cheaper check.  In our binding, the producer is modelled by an unrestricted
envelope and the consumer is Type-3; the verifier-side conformance
budget falls to $O(n)$ per candidate trace by
Theorem~\ref{thm:closure}.

A natural question follows.  \emph{May a Type-3 agent itself act
as a producer?}  The answer is yes, with one essential condition.
Concretely, in our system the verifier may discover a new
fixture (e.g., a new column-name keyword as the train CSV schema
is augmented mid-run) and \emph{push} that keyword into the
keyword set $K$.  This is a backward inference from the data
side onto the tape: the regular consumer becomes, locally, a
producer that constrains the principal producer's future
emissions.

The condition is decidability.  The Type-3 producer's loop must
terminate in a decidable verdict, otherwise the classical
halting problem of Turing~\cite{turing1936computable}
re-asserts itself on the consumer side and the architecture
loses its verifier-side budget.  Concretely:  if the
keyword-update rule is itself Type-0 (e.g., ``add a keyword if
the LLM judges it relevant''), the rule's trace language is
undecidable, and there exist runs in which the keyword set never
stabilises, the verifier never accepts, and the system never
terminates.  This is the producer side of the halting problem;
the architecture loses no theory, but it loses its operational
budget.

The escape is to make the keyword-update rule itself Type-3
--- i.e., a regular grammar over a finite catalogue of
admissible keyword extensions --- or, where this is infeasible,
to insert a final \emph{human} verification and validation step
that closes the loop.  In the realised system, both escapes are
used.  The \textsf{test\_*} extensions are admitted by a regular
grammar read off the train CSV schema; the leaderboard
submission itself is reviewed by the first author before upload.
The system is therefore an idealised producer constrained by a Type-3 verifier in the
inner loop, and a human-in-the-loop closure on the outer loop.
Both are sound; neither, alone, would be.

This is the operational meaning of ``yes, with software
verification and validation techniques inferred into the
Data-Driven agent.''  The V\&V techniques are precisely the
machinery of Hoare~\cite{hoare1969axiomatic},
Dijkstra~\cite{dijkstra1975guarded},
Pnueli~\cite{pnueli1977temporal}, and
Clarke--Emerson~\cite{clarke1981ctl,clarke1999modelchecking}
applied at the regular level of the Aho--Corasick consumer.
They are not a metaphor; they are the keyword set $K$, the DFA
$M_K$, and the LTL formula
$G(\mathit{plan} \rightarrow F\,\mathit{verified})$.

\subsection{Threats to validity}

\paragraph{Construct validity.}
The trace-language identity criterion (Section~\ref{sec:trace})
is a strong claim: an agent is \emph{nothing} but its trace
language.  The criterion is well-defined --- two agents are
identified iff $L(\mathcal{A}_1) = L(\mathcal{A}_2)$ --- but it
does discard information about, e.g., the LLM's chain-of-thought
or the planner's intermediate JSON.  We accept the discard
because the discarded objects are not directly verifiable; only
the trace is, and only the trace is what we admit as evidence.

\paragraph{Internal validity.}
The Goal-Driven producer's LLM core is non-deterministic.  Two
re-runs of Algorithm~\ref{alg:enum} on the same input produce
slightly different traces.  The class signature
(Table~\ref{tab:class-signature}) is robust across re-runs in
our pilot experiments, but the leaderboard placement varies by
$\pm$10--20 ranks across re-runs.  We report a single placement
($710/2330$) and acknowledge the residual variance.  The
reproducibility kit at\\
\url{https://www.kaggle.com/sweeden-ttu/competitions/spaceship-titanic}\\
allows external re-runs at fixed random seeds.

\paragraph{External validity.}
The architecture has been instantiated on one task
(\textsc{Spaceship Titanic}).  We expect the trace-language
account to generalise to other Kaggle-style tabular tasks
(\textsc{House Prices}, \textsc{Titanic},
\textsc{Digit Recognizer}), because the start language and
consumer are task-agnostic.  We do not claim generalisation
outside the tabular regime; agents that operate on images, text,
or open-ended code may require an expanded start language and a
correspondingly expanded keyword set $K$.

\paragraph{Coverage of the producer--consumer escape.}
The discussion in Section~\ref{sec:type3-producer} relies on
either (a) a regular keyword-update rule, or (b) a human in the
loop.  In our experiments both held; we do not claim a general
solution when neither holds.  The recursion-without-decidability
case is exactly where the halting problem returns and the
architecture provides no warranty.

\subsection{Implications for the field}

Two implications follow.  First, the formal-language framing of
agents need not be metaphorical.  When done at the level of
narrowness of this paper, it produces concrete bounds on
verifier-side conformance cost (Theorem~\ref{thm:closure},
Corollary~\ref{cor:budget}) and operational architectures (one
idealised producer, one Type-3 consumer) that are competitive with
the contemporary state of the art.

Second, the multiplication of agent ``roles'' in contemporary
literature --- planner, executor, critic, verifier, summariser
--- is, at least in our case, an emergent phenomenon, not an
architectural choice.  The architect designs two agents; the
binding produces five more.  This is consistent with
McCarthy's~\cite{mccarthy1960lisp} programs-as-data and
Newell--Simon's~\cite{newell1976physical} symbol-system
hypothesis: roles are projections, not primitives.  We
conjecture, but do not here prove, that other contemporary agent
architectures admit a similar derivation: a small number of
principal agents binds a verifier, and the named roles are
projections of the validated trace.

\subsection{What we did not do}

We did not measure the LLM's chain-of-thought.  We did not
compare backbone families.  Most importantly, we did not yet run the
ablation that would most directly test the theory: no verifier versus
Aho--Corasick verifier versus LLM-as-verifier, measured by acceptance
rate, trace length, runtime, and final leaderboard score.  That
experiment is the clearest next test of the claim that a Type-3
verifier meaningfully constrains a more expressive generator.  We
also did not formally prove that the realised trace is the unique
fixed point of the binding.  Each is a defensible piece of follow-up
work; the present claim is narrower: trace languages provide a useful
specification lens, and a regular verifier gives a linear-time
conformance check that can guide a contemporary ML-engineering
workflow.
